\chapter[Was ich war und was ich jetzt bin]{Was ich war und was ich jetzt bin\footnote{Mit der Bibel in der Hand zu lesen.}}
\quellseite{10}

Lieber Freund! Ich wohnte einst bei 2 Tim. 3, 4. und wandelte in Eph. 2, 2., und meine tägliche Lebensweise steht noch beschrieben Eph. 2, 3. Da hörte ich eines Tages, daß ein Erbgut für mich vorhanden sei, und erhielt auch eine Beschreibung davon, die du noch 1 Petr. 1, 4. lesen kannst. Ein gütiger Herr, der in Hebr. 4, 14. wohnt, hatte es mir erworben und den Preis dafür mit 1 Petr. 1, 19. bezahlt. Aber, die Wahrheit zu gestehen, ich glaubte der Angabe nicht, da ich den Mann gar nicht kannte und aus langjähriger Erfahrung wußte, daß wir von Fremden keine Gunstbezeugungen aus bloßer Liebe zu erwarten hatten, und von Freunden selten solche, die viel kosten. Dennoch sprach ich in 2 Tim. 3, 16. vor, da meine Aussichten in Eph. 2, 12. so schlecht waren, wie sie nur sein konnten.

Ich fand denn das Haus, das ich suchte, in 2 Cor. 5, 1., und die Bedingungen, unter denen es zu haben war — du findest sie aufgezeichnet Jes. 55, 1. 2. und Joh. 7, 37. —, waren in der That wunderbar einladend für Arme und Hülfsbedürftige. Das Haus hatte aber nur eine einzige Thür, und es dauerte eine Zeitlang, ehe ich sie in Joh. 10, 9. fand. Meine bleibende Wohnung wird nun in 2 Cor. 5, 1. sein. Aber wenn du in Hebr. 4, 16. einmal vorsprechen willst, wirst du mich und viele andere, meine Brüder und Schwestern, dort finden, weil wir dort täglich zusammentreffen. Wenn du kommen willst, so kannst du dich auf das, was der Knecht Luc. 14, 22. sagt, fest verlassen.
